\documentclass[12pt]{article}

%% Basic document formatting
\usepackage{amsmath, amsthm, amssymb, amsfonts}
\usepackage{mathtools}
\usepackage{mathrsfs}
\usepackage{xspace}
\usepackage{thmtools}
\usepackage{graphicx}
\usepackage{tikz-cd}
\usepackage{color}
\usepackage{setspace}
\usepackage{fancyhdr}
\usepackage{titling}
\usepackage{hyperref}
\usepackage[left=0.4in,right=0.4in,top=1in,bottom=1in]{geometry}
\usepackage{float}
\usepackage{tabularx}
\usepackage{booktabs}
\usepackage[utf8]{inputenc}
\usepackage[english]{babel}
\usepackage{framed}
\usepackage[dvipsnames]{xcolor}
\usepackage{environ}
\usepackage{tcolorbox}
\tcbuselibrary{theorems,skins,breakable}

\usepackage{enumitem}
\setlist[enumerate]{leftmargin=*}
\setlist[enumerate,1]{labelindent=\parindent}
\setlist[enumerate,2]{labelindent=0pt}

% Blackboard Bold
\newcommand{\Z}{\mathbb{Z}}                 % integers
\newcommand{\Zpl}{\mathbb{Z}_{+}}           % positive integers
\newcommand{\Znn}{\mathbb{Z}_{\geq 0}}      % nonnegative integers. 
\newcommand{\Q}{\mathbb{Q}}                 % rationals
\newcommand{\Qpl}{\mathbb{Q}_{+}}           % positive Rationals
\newcommand{\R}{\mathbb{R}}                 % reals
\newcommand{\Rpl}{\mathbb{R}_{+}}           % positive Reals
\newcommand{\C}{\mathbb{C}}                 % complex numbers
\newcommand{\F}{\mathbb{F}}                 % field

% Words
\newcommand{\st}{\text{ such that }}
\newcommand{\wrt}{\text{ with respect to }}
\newcommand{\with}{\text{ with }}
\newcommand{\ie}{\text{, i.e. }}
\newcommand*{\ditto}{\text{---}\texttt{"}\text{---}}

% Operators
\newcommand{\abs}[1]{\left|#1\right|}                   % absolute value
\newcommand{\norm}[1]{\abs{\abs{#1}}}
\newcommand{\floor}[1]{\left\lfloor #1 \right\rfloor}   % floor
\newcommand{\ceil}[1]{\left\lceil #1 \right\rceil}      % ceiling
\newcommand{\powset}[1]{\mathscr{P}(#1)}

% Category Theory 
\renewcommand{\hom}{\text{Hom}}
\newcommand{\homspace}[3]{\hom_{#1}(#2, #3)}
\newcommand{\coproduct}[9]{
  \begin{tikzcd}[ampersand replacement=\&]
      \& \& #1 \arrow[dll, bend right, "#6"] \arrow[dl, "#5"] \\
   #4 \& #3 \arrow[l, "#9"] \\
      \& \& #2 \arrow[ull, bend left, "#8"] \arrow[ul, "#7"]
  \end{tikzcd}
}

\newcommand{\product}[9]{ 
  \begin{tikzcd}[ampersand replacement=\&]
      \& \& #3 \\
      #1 \arrow[r, "#7"] \arrow[urr, bend left, "#5"] \arrow[drr, bend right, "#6"] \& #2 \arrow[ur, "#8"] \arrow[dr, "#9"] \\ 
      \& \& #4
  \end{tikzcd}
}


% Algebra 
\newcommand{\Syl}[2]{\text{Syl}_{#1}{(#2)}}           % set of Sylow-p subgroups 
\newcommand{\<}{\langle}                % \<x\>, subgroup generated by x 
\renewcommand{\>}{\rangle}
\renewcommand{\index}[2]{[#1 : #2]}
\newcommand{\id}{\text{id}}             % identity element
\newcommand{\lhdneq}{%
  \mathrel{\ooalign{$\lneq$\cr\raise.22ex\hbox{$\lhd$}\cr}}}
\newcommand{\order}[1]{\text{o}(#1)}    % order of an element
\newcommand{\spec}{\operatorname{Spec}}
\newcommand{\rad}{\operatorname{rad}}   % radical of an ideal 
\newcommand{\sym}[1]{\mathfrak{S}_{#1}}
\newcommand{\aut}[1]{\text{Aut}(#1)} 
\newcommand{\gal}[2]{\text{Gal}(#1 / #2)} 
\newcommand{\inn}[1]{\text{Inn}(#1)} 
\let\oldcong\cong
\let\oldequiv\equiv
\renewcommand{\cong}{\oldequiv}
\renewcommand{\equiv}{\oldcong}

\newcommand{\triangleleftneq}{\mathrel{\ooalign{%
  $\trianglelefteq$\cr
  \hidewidth\raise0.06ex\hbox{$\not\mkern4mu$}\hidewidth\cr
}}}

\makeatletter % cycle
\newcommand{\cyc}[1]{(\mathbf{\cyc@process#1\relax})}
\def\cyc@process#1#2\relax{%
	#1%
	\ifx\relax#2\relax
	\else
		\,\cyc@process#2\relax
	\fi
}
\makeatother

\newcommand{\GL}[2]{\text{GL}_{#1}(#2)}                             % general linear group
\newcommand{\SL}[2]{\text{SL}_{#1}(#2)}                             % special linear group
\newcommand{\gl}[2]{\mathfrak{gl}_{#1}(#2)}
\renewcommand{\sl}[2]{\mathfrak{sl}_{#1}(#2)}
\newcommand{\so}[2]{\mathfrak{so}_{#1}(#2)}
\newcommand{\su}[1]{\mathfrak{su}(#1)}

% Linear Algebra
\newcommand{\bmat}[1]{\begin{bmatrix}#1\end{bmatrix}}   % bracketed matrix
\newcommand{\rank}{\operatorname{rank}}                 % rank
\newcommand{\nullity}{\operatorname{nullity}}           % nullity
\newcommand{\tr}{\operatorname{tr}}

% Combinatorics
\newcommand{\qnum}[1]{\left[#1\right]_q}                % q-analog
\newcommand{\qfact}[1]{\left[#1\right]!_q}              % q-analog factorial 
\newcommand{\qbinom}[2]{\binom{#1}{#2}_q}               % Gaussian binomial coefficient

% Topology/Analysis
\newcommand{\ball}[2]{\text{B}_{#1}(#2)}  % B_r(x): open r-balls around x
\newcommand{\diam}{\text{diam}}           % diamter of a set in metric space 
\newcommand{\lowint}[2]{\underline{\int_{#1}^{#2}}}
\newcommand{\upint}[2]{\overline{\int}_{#1}^{#2}}

% Misc Notation
\newcommand{\oneton}[1]{\{1, \ldots, #1\}}
\newcommand{\defeq}{\vcentcolon=}   % :=
\newcommand{\eqdef}{=\vcentcolon}   % =: 
\renewcommand{\bf}[1]{\textbf{#1}}
\newcommand{\im}{\operatorname{im}}
\newcommand{\fnrest}[2]{\left.#1\right|_{#2}}
\newcommand{\isom}{\overset{\sim}{\rightarrow}} 


\renewcommand{\thesection}{\arabic{section}.}
\renewcommand{\thesubsection}{(\alph{subsection})}

\newtheorem{claim}{Claim}
\newtheorem*{lemma}{Lemma}

%% Headers & title setup
\newcommand{\course}{Math100C}
\newcommand{\myname}{Jay Ser}
\setlength{\headheight}{14.5pt}
\pagestyle{fancy}
\fancyhf{}
\renewcommand{\headrulewidth}{0.4pt}
\lhead{\course}
\rhead{\myname}
\cfoot{\thepage}
\setlength{\droptitle}{-4em} 
\title{\course\ - WK8} 
\author{\myname}
\date{2026.05.22}

\begin{document}
\maketitle
\thispagestyle{fancy}

\newtheorem*{hint}{Hint}

%------------------------------------------------------------------------------%

\section{} 
Let $n$ be a positive integer. 
Define the polynomials 
$$s_i = u_1^i + \ldots + u_n^{i}, \quad (i \in \oneton{n})$$ 
Prove that these generate the ring of symmetric polynomials in $u_1, \ldots, u_n$ over $\Q$ but not over $\Z$ when $n \geq 2$. 

\noindent\dotfill 

Newtons' identities (plural?) gives a recursive formula for calculating $ke_k$, $k \in \oneton{n}$ from $s_1, \ldots, s_k$ and $e_1, \ldots, e_{k - 1}$: 
\begin{equation} \label{eq:1_newton} 
  ke_k = \sum_{i = 1}^{k} e_{k - i} s_{i}
\end{equation}
Over $\Q$, the coefficient $k$ on the left hand side of \eqref{eq:1_newton} can be canceled, so Newton's identities says all the elementary symmetric polynomials $e_1, \ldots, e_n$ can be written in terms of $s_1, \ldots, s_n$; 
the elementary symmetric polynomials generate the ring of symmetric polynomials over coefficients in any ring $R$, so $s_1, \ldots, s_n$ generate the ring of symmetric polynomials over $\Q$. 

I don't really know how to proceed for the the symmetric polynomials over $\Z$. 

\section{} 
Let $n$ be a positive integer and $G$ a subgroup of $\mathfrak{S}_n$ containing at least one transposition. 
\subsection{} 
Suppose $G$ contains all the transpositions of $\mathfrak{S}_n$. 
Prove $G = \mathfrak{S}_n$. 

\noindent \dotfill

Since every permutation in $\sym{n}$ has a unique cycle decomposition, it suffices to show that $G$ contains all the cycles of $\sym{n}$. 
To do this, induct on cycle length $k$.
$k = 2$ holds by the assumption on $G$. 
Suppose $G$ contains all cycles of length less than $k$. 
Since relabeling the elements $1, 2, \ldots, n$ doesn't affect the inductive hypothesis, showing that $G$ contains all $k$-cycles is equivalent to showing that $G$ simply contains $\cyc{12 \cdots k}$. 
Well, by the inductive hypothesis, $\cyc{12 \cdots [k - 1]}, \cyc{[k - 1]k} \in G$, so 
$$ \cyc{12 \cdots [k - 1]} \cdot \cyc{[k - 1]k} = \cyc{12 \cdots k} \in G,$$
finishing the proof. 

\subsection{} 
Suppose $n = p$ is prime and $p \mid \#G$. 
Prove $G = \mathfrak{S}_n$ by reducing to the previous question. 

\noindent\dotfill 

I first prove the lemma Kiran suggested: 
\begin{lemma}
  Let a subgroup $G < \sym{n}$ correspond to a graph $\Gamma = (V, E)$ with vertices $1, \ldots, n$ and edge $ij$ whenever $G$ contains a transposition $\cyc{ij}$. 
  If $\Gamma$ is connected, then $G = \sym{n}$. 
\end{lemma} 
\begin{proof}[Proof of Lemma] 
  First, consider the special case where the vertex 1 is connected to all other vertices, i.e., 
  $G$ contains all transpositions $\cyc{1k}$ for $k \in \{2, \ldots, n\}$. 
  Then $G$ contains all the transpositions $\sym{n}$: indeed, for any $i \neq j \in \oneton{n}$, 
  $$\cyc{1 i}\cyc{1 j}\cyc{1 i} = \cyc{ij}$$
  By part (a), $G = \sym{n}$. 

  The general case will reduce to this special case. 
  Given any $k \in \{2, \ldots, n\}$, $\cyc{1k} \in G$ by the following: 
  since $\Gamma$ is connected, there is a walk 
  $$1 \rightsquigarrow u_1 \rightsquigarrow \ldots \rightsquigarrow u_t = k$$
  But the calculation 
  $$\cyc{1 {u_1}} \cyc{{u_1} {u_2}} \cyc{1 {u_1}} = \cyc{1 {u_2}}$$ 
  shows there is a direct path $1 \rightsquigarrow u_2$ in $\Gamma$. 
  Inductively performing this calculation, one simply obtains a direct path $1 \rightsquigarrow k$, i.e., 
  $\cyc{1 k} \in G$, as was to be shown. 
\end{proof} 

Let $G$ be a subgroup of $\sym{p}$ such that it contains a transposition and $p \mid \#G$. 
By Cauchy's theorem, $G$ contains an element of $\sym{p}$ of order $p$. 
The only such element is a $p$-cycle; 
by relabeling the $p$-cycle as needed, one can assume that $G$ contains $\cyc{12 \cdots p}$ and some transposition $\cyc{ij}$. 
$$\cyc{12 \cdots p} \cyc{ij} \cyc{1 p {[p - 1]} \cdots 2} = \cyc{{[i + 1]}{[j + 1]}}$$
shows $G$ contains all transpositions of the form $\cyc{{[i + t]}{[j + t]}}$. 
Namely, it contains $\cyc{1k}$ for some $k \in \{2, \ldots, p\}$ and therefore $\cyc{{[1 + t]} {[k + t]}}$ for all $t$. 
If $\Gamma$ is the graph corresponding to $G$ as in the lemma, the notice that $\Gamma$ has vertices 
$$1 \rightsquigarrow k \rightsquigarrow (2k - 1) \rightsquigarrow (3k - 2) \rightsquigarrow \ldots$$
Hence if 
$$x \mapsto xk - (x - 1) = (k - 1)x + 1$$
is surjective mod $p$, then $\Gamma$ is connected. 
Indeed, the map is surjective mod $p$ because $k - 1 \neq 0$ mod $p$ and linear equations always have a solution mod $p$ (since $\F_p$ is a field). 
Hence $\Gamma$ is connected, and the lemma then gives $\sym{p} = G$. 

\section{}
Let $f$ be a polynomial of degree $n$ over a field $F$. 
Prove that a splitting field for $f$ over $F$ has degree dividing $n!$

\noindent\dotfill

Let $K / F$ be the splitting field for the polynomial $f$ over $F$ of degree $n$. 
Letting $G \defeq \gal{K}{F}$, recall $K / F$ is a Galois extension, $K^{G} = F$, and $[K : K^{G}] = \#G$.  
Since $K / F$ is Galois, $G$ acts on the roots of $f$ faithfully, giving
$$G \hookrightarrow \mathfrak{S}_n$$ 
hence $\#G \mid \# \mathfrak{S}_n = n!$. 
Putting together all the facts, 
$$[K : F] \mid n!,$$
as was to be shown. 

\section{}
Compute the Galois group of the splitting field $x^5 - 2$ over $\Q$. 

\noindent\dotfill 

The splitting field for $f(x) \defeq x^5 - 2$ over $\Q$ is 
$$\Q(\zeta_5, 2^{1/5}):$$ 
$\zeta_5$ and $2^{1/5}$ clearly generate all the roots of $f$; 
meanwhile, any field containing all the roots of $f$ clearly contains $2^{1/5}$ and $\zeta_{5} 2^{1/5}$, and multiplying the inverse of the former with the latter gives $\zeta_5$ in that field as well. 

Let $K \defeq \Q(\zeta_5, 2^{1/5})$. 
The minimal polynomial for $2^{1/5}$ over $\Q$ is $f$, a degree 5 polynomial, whereas $\zeta_5$ is a degree 4 element over $\Q$. 
Since these are coprime,
$$\index{K}{\Q} = 20.$$
Finally, it is clear $f$ is irreducible, for example by Eisenstein for $p = 2$, so the Galois group $G \defeq \gal{K}{\Q}$ is the unique transitive subgroup of $\sym{5}$ of order 20. 

Let's compute this group. 
Associate the roots of $f$ $2^{1/5}, 2^{1/5} \zeta_5, \ldots, 2^{1/5} \zeta_5^{4}$ with indices $1, 2, \ldots, 5$, respectively. 
First, consider the permutation $\sigma \in G$ mapping 
$$\zeta_5 \mapsto \zeta_5^{2}, \quad 2^{1/5} \mapsto 2^{1/5}.$$ 
One can easily check that $\sigma$ has order 4 in $G$ and acts on the roots of $f$ as $\cyc{2354}$. 
Next, consider $\tau \in G$ mapping 
$$\zeta_5 \mapsto \zeta_5, \quad 2^{1/5} \mapsto 2^{1/5} \zeta_5.$$
It is clear that $\tau$ is order 5 in $G$ acting on the roots of $f$ as $\cyc{12345}$. 
A 5-cycle and a 4-cycle generates a group of order 20, i.e., 
$$G = \<\sigma, \tau\>.$$
This computes the Galois group for the splitting field of $x^5 - 2$ over $\Q$. 

\section{}
Prove that $\Q(\sqrt{2}, \sqrt{3}, \sqrt{5})$ is a Galois extension of $\Q$ and compute its Galois group. 

\noindent \dotfill

Let $K \defeq \Q(\sqrt{2}, \sqrt{3}, \sqrt{5})$. 
The generators of $K$ over $\Q$ are precisely the roots (up to $\pm$) of 
$$f(x) \defeq (x^2 - 2)(x^2 - 3)(x^2 - 5),$$
so $K$ is a splitting field over $\Q$, 
i.e., $K / \Q$ is a Galois extension. 
We've shown many times in class that $\Q(\sqrt{2}, \sqrt{3})$ is a degree 4 extension of $\Q$, and it's clear $\sqrt{5} \notin \Q(\sqrt{2}, \sqrt{3})$, so 
$$\index{K}{\Q} = 8.$$

Let $G \defeq \gal{K}{\Q}$. 
By the paragraph above, $G$ is a subgroup of $\sym{6}$ of order 8. 
Since $K$ is the splitting field for $f$, any $\sigma \in G$ permutes the roots of each factor of $f$: 
$$\sqrt{2} \mapsto \pm \sqrt{2}, \quad \sqrt{3} \mapsto \pm \sqrt{3}, \quad  \sqrt{5} \mapsto \pm \sqrt{5}$$
None of the factors of $f$ share a root, so each of these three choices can be made independently; 
i.e., the permutations of $G$ sending 
\begin{itemize}
  \item $\sqrt{2} \mapsto -\sqrt{2}$, identity on other elements 
  \item $\sqrt{3} \mapsto -\sqrt{3}$, idnetity on other elements 
  \item $\sqrt{5} \mapsto -\sqrt{5}$, identity on other elements 
\end{itemize}
correpond to disjoint $2$-cycles, hence they generate $(\Z / 2\Z) \times (\Z / 2\Z) \times (\Z / 2\Z)$, which has order 8. 
Hence $G \equiv (\Z / 2\Z)^3$.

\section{}
Define a field $K = \C(t)$. 
Define automorphisms $\sigma, \tau$ of $K$ over $\C$ mapping 
$$\sigma(t) = \frac{t + i}{t - i}, \quad \tau(t) = \frac{it - i}{t + 1}$$
Prove these automorphisms generate a group isomorphic to $A_4$ and compute the fixed field. 

\noindent \dotfill 

I utilize the isomorphism $\GL{2}{\C} / \C \equiv \text{Aut}(\C(t) / \C)$. 
Identify 
\begin{gather*}
  \sigma(t) \rightsquigarrow M_\sigma \defeq \bmat{1 & i \\ 1 & -i} \\ 
  \tau(t) \rightsquigarrow M_\tau \defeq \bmat{i & -i \\ 1 & 1}. 
\end{gather*}
Then 
\begin{gather*}
  (M_\sigma)^{2} = \bmat{1 + i & 1 + i \\ 1 - i & i - 1}, \quad (M_\sigma)^{3} = \bmat{2 + 2i & 0 \\ 0 & 2 + 2i} = I_2 \\ 
  (M_\tau)^{2} = \bmat{-1 - i & 1 - i \\ i + 1 & 1 - i}, \quad (M_\tau)^{3} = \bmat{2 -2i & 0 \\ 0 & 2 - 2i} = I_2
\end{gather*}
So $(M_\sigma)^{3} = (M_\tau)^{3} = 1$. 
Also, 
$$M_\sigma M_\tau = \bmat{2i & 0 \\ 0 & -2i} = \bmat{1 & 0 \\ 0 & -1}$$
shows $M_\sigma M_\tau \neq I_2$ and $(M_\sigma M_\tau)^{2} = I_2$. 
This means $M_\sigma$ and $M_\tau$ satisfy the relations for the presentation 
$$A_4 = \<x, y \mid x^3 = y^3 = (xy)^2 = 1\>$$
In other words, $\sigma$ and $\tau$ generate $A_4$. 

I'm unsure what the fixed field is. 
I suspect it is some sum of the matrices computed above. 

\section{} 
Let $\zeta_n \defeq e^{2\pi / n}$ for some positive integer $n$. 
Prove $\Q(\zeta_n)$ is a Galois extension of $\Q$ with Galois group isomorphic to $(\Z / n\Z)^{\times}$. 

\noindent \dotfill

Let $G = \gal{\Q(\zeta_n)}{\Q}$. 
Recall that the minimal polynomial for $\zeta_n$ over $\Q$ is 

\begin{equation} \label{eq:7_cyc} 
  \Phi_n(x) = \prod_{\substack{a < n \\ (a, n) = 1}} (x - \zeta_n^{a})
\end{equation}
Thus it is clear that $\Q(\zeta_n) / \Q$ is a splitting field for $\Phi_n(x)$ over $\Q$,
i.e., a Galois extension of $\Q$: 
it contains all the roots of $\Phi_n(x)$, while obviously, any field containing all the roots of $\Phi_n(x)$ contains $\zeta_n$. 
\eqref{eq:7_cyc} also shows that any $\sigma \in G$ maps 
$$\zeta_n \mapsto \zeta_n^{a}$$
for some $a < n$, $(a, n) = 1$. 
Hence the following map $G \rightarrow (\Z / n\Z)^{\times}$ is well defined: 
\begin{equation} \label{eq:7_isom} 
  \sigma \mapsto a 
\end{equation}
where $a$ is such that $\sigma(\zeta_n) = \zeta_{n}^{a}$. 
This map is clearly injective since $\sigma$ is defined entirely by how it sends $\zeta_n$, which means only one $\sigma \in G$ sends $\zeta_n$ to $\zeta_n^{a}$ for a fixed $a$. 
The map is also a homomorphism: 
given $\sigma_1, \sigma_2 \in G$ with $\sigma_1(\zeta_n) = \zeta_n^{a_1}$ and $\sigma_2(\zeta_n) = \zeta_n^{a_2}$, 
\begin{align*}
  \sigma_1 \sigma_2 (\zeta_n) & = \sigma_1 (\zeta_n^{a_2}) \\ 
							  & = (\sigma_1(\zeta_n))^{a_2} \\ 
							  & = \zeta_n^{a_1 a_2} 
\end{align*}
Finally, because $\Q(\zeta_n) / \Q$ is a splitting field, $\#G = \varphi(n)$. 
$\varphi(n) = (\Z / n\Z)^{\times}$ as well, so \eqref{eq:7_isom} is surjective; 
this finishes the proof that $G \equiv (\Z / n\Z)^{\times}$. 

\section{}
Let $f(x) = x^3 + x + 1$. 
Prove the splitting field of $f$ over $\Q$ contains $\sqrt{-31}$. 

\noindent\dotfill 

I read about the discriminant not in Artin, but in Dummit and Foote because the latter makes more sense to me. 
In general, for any $f$, the definition of the discriminant gives 
$$\sqrt{D} = \prod_{i < j}(\alpha_i - \alpha_j)$$
where the $\alpha_i$ and $\alpha_j$'s are roots of $f$, which shows $\sqrt{D}$ is always in the splitting field for $f$. 
For $f(x) = x^3 + x + 1$, the discriminant is 
$$D \defeq -4(1)^3 - 27(1)^2 = -31.$$
Hence $\sqrt{-31}$ is in the splitting field for $f$ over $\Q$. 

\end{document}
